



\section{Introduction}
\label{sec:intro}




When constructing a transformer-based language model, a major design decision is the length of training sequences, denoted $L$ herein, which has to date  been equivalent to the length of inference sequences.  More context, achieved by larger $L$, improves predictions at inference time.  But longer sequences are more expensive to train on.\footnote{Figure~\ref{fig:wt103_train_speeds_all} in the appendix plots training speed, in words per second, against $L$.}  


Before transformers, RNN language models were trained on shorter-$L$ sequences and assumed to generalize to longer contexts at inference time~\citep{Mikolov2010RecurrentNN, Mikolov2012ContextDR, zaremba2014recurrent}.  \citet{vaswani}, introducing the transformer, speculated that it ``may [...] extrapolate to sequence lengths longer than the ones encountered during training.''  
We define \emph{extrapolation} as a model's ability to continue performing well as the number of input tokens during validation increases beyond the number of tokens on which the the model was trained. 
We find that transformer language models (LMs)  that use sinusoidal position embeddings have very weak  %
extrapolation abilities; see Figure~\ref{fig:wt103_extra}.



\begin{figure}[h]
\centering
\begin{subfigure}{.5\textwidth}
  \centering
  \includegraphics[width=\linewidth]{figures/wt103_extra_512.pdf}
\end{subfigure}%
\begin{subfigure}{.5\textwidth}
  \centering
  \includegraphics[width=\linewidth]{figures/wt103_extra_1024.pdf}
\end{subfigure}
\caption{Extrapolation:
as the (validation-set's) input sequence gets longer ($x$-axis), current position methods (sinusoidal, rotary, and T5) show degraded perplexity ($y$-axis, lower is better), but our method (\S\ref{sec:ourmethod}) does not.  Models were trained on WikiText-103 with sequences of $L=$ 512 (left) or $L=$ 1,024 (right) tokens.  T5 ran out of memory on our 32GB GPU. For more detail on exact perplexities and runtimes, see Tables~\ref{tab:appendix:wt103_extra_512} and~\ref{tab:appendix:wt103_extra_1024} in the appendix.}

\label{fig:wt103_extra}
\end{figure}











We demonstrate that this failure to extrapolate is caused by the position embedding method.  As shown in Figure~\ref{fig:wt103_extra}, recent alternatives to the original sinusoidal position method \citep{roformer,t5} have improved extrapolation.  However, the better of these, the T5 bias, is considerably slower than the sinusoidal approach and uses extra memory and parameters (Figure~\ref{fig:wt103_speed_mem}).



We therefore introduce Attention with Linear Biases (ALiBi) to facilitate  efficient extrapolation. %
ALiBi negatively biases attention scores with a linearly decreasing penalty proportional to the distance between the relevant key and query. Our simple approach eliminates position embeddings. %



Compared to a sinusoidal model trained on the same input length, our method requires no additional runtime or parameters and incurs a negligible (0--0.7\%) memory increase.  ALiBi can be implemented by changing only a few lines of existing transformer code. 


Using ALiBi, a transformer LM can be trained on short-$L$ sequences and therefore at much lower cost, and it can still be reliably applied to long sequences at runtime. For example, a 1.3 billion parameter LM trained on $L=$ 1024 tokens with ALiBi achieves the same perplexity as a sinusoidal model trained on $L=$ 2048 when both are tested on sequences of 2048 tokens, even though \textit{our model is 11\% faster and uses 11\% less memory. }%


Though performance peaks at around two times the number of tokens that the model was trained on, ALiBi maintains strong performance even on sequences of length 10,000. 
In recently explored settings where NLP training examples are given as context to an LM \citep{gpt3}, our approach will allow exposure to more examples. Additionally, it enables generation of longer outputs.












